\section{Conclusiones}

Los resultados obtenidos permitieron analizar los principales factores que influyen en el movimiento de un péndulo simple y comparar su comportamiento con el modelo teórico. Para las tres masas ensayadas se obtuvieron períodos de $2.08\,\mathrm{s}$, $2.05\,\mathrm{s}$ y $2.07\,\mathrm{s}$, respectivamente, con desviaciones pequeñas respecto al valor teórico calculado para la longitud utilizada. Estas variaciones no mostraron una tendencia sistemática con la masa, por lo que los resultados fueron consistentes con la predicción del modelo de péndulo simple, según la cual el período no depende de la masa.

Al modificar la longitud del péndulo entre $99.5\,\mathrm{cm}$ y $79.5\,\mathrm{cm}$, se observó una tendencia general a la disminución del período, en concordancia con la relación teórica $T\propto\sqrt{L}$. Aunque algunas configuraciones presentaron pequeñas desviaciones respecto a la tendencia ideal, estas pueden asociarse a las limitaciones del montaje y a las incertidumbres propias del procedimiento experimental.

Para las amplitudes angulares comprendidas entre $15^\circ$ y $30^\circ$, el período permaneció aproximadamente constante, con una diferencia relativa máxima de \SI{1.34}{\percent} entre los ensayos. Este comportamiento fue consistente con la aproximación de pequeñas oscilaciones utilizada en el modelo del péndulo simple.

En el estudio del régimen amortiguado se observó una disminución progresiva de la amplitud con el tiempo. El ajuste exponencial aplicado a los máximos de amplitud presentó un coeficiente de determinación $R^2=0.9693$, lo que indicó una buena correspondencia entre el modelo exponencial y los datos experimentales. A partir de este ajuste se obtuvo una constante de amortiguamiento $b=(6.65\pm0.32)\times10^{-5}\,\mathrm{kg/s}$. El comportamiento observado fue compatible con un régimen subamortiguado y con el modelo de amortiguamiento viscoso lineal utilizado para describir la disipación de energía del sistema.

En conjunto, los resultados mostraron concordancia general con las predicciones del modelo de péndulo simple, aunque el montaje real presentó desviaciones respecto a las condiciones ideales que pueden estar asociadas con el rozamiento, la resistencia del aire, la masa y elasticidad del cordel, la determinación de la longitud efectiva y la liberación manual del péndulo.

Como posibles mejoras experimentales, se podría emplear un soporte más rígido, reducir las variaciones asociadas a la liberación manual y utilizar esferas de menor tamaño para disminuir la incertidumbre en la determinación de su centro geométrico. Como extensión del experimento, podría estudiarse el comportamiento para amplitudes mayores que $30^\circ$, con el objetivo de analizar la pérdida de validez de la aproximación $\sin\theta\approx\theta$, o comparar el amortiguamiento del sistema bajo diferentes condiciones experimentales.
